% Pursuing Stacks (cotes 134-2 to 134-8): archive cote 134-3, site page 51 (PDF page 52), typescript page 51.
% typescript. Draft compared: ps2.tex:2550-2624.
% AI-transcribed by Claude Opus 5.5 (claude-opus-5-5, adaptive thinking, effort high), Message Batches API, batch ps1 (msgbatch_01Rv5hCaNTtUt4JHx1yqoHPL), 2026-11-12; not yet checked by eye.
\scanpage{134-3/51}
\mspage{51}
X, the cohomology \xout{groups} $H^i(X,G)$ \inserted{$(i\hand{\geq}1)$,} is trivial whenever defined, namely \deleted{any}\inserted{for}
$i\hand{\geq}1$ if G commutative, and $i=1$ if G is not supposed commut\xout{t}ative. The
0-connexity is readily translated into a corresponding property of the fi-%
nal element in the category of sheaves $Sh(X)=\hand{\mathcal{A}}$ on X. An element \xout{\unsure{of}} U of $\hand{\mathcal{A}}$
is called \uline{aspheric} if the induced topos X/U or $\hand{\mathcal{A}}_{/U}$ is aspheric. An arrow
$f\colon U'\hand{\to} U$ in $\hand{\mathcal{A}}$ is called a \uline{weak equivalence} resp. \uline{aspheric}, if the cor-%
responding ``map'' or morphism of the induced topoï
\[ X/U'\hand{\longrightarrow} X/U \quad\text{or}\quad \hand{\mathcal{A}}_{/U'}\hand{\longrightarrow}\hand{\mathcal{A}}_{/U} \]
has the corresponding property. Of course, in the notations it is often
convenient simply to identify objects \inserted{U} of $\hand{\mathcal{A}}$ with the induced topos, and
accordingly for arrows. \xout{The \unsure{notion of}} Clearly if f is aspheric it is a
weak equivalence, more specifically, f is \xout{a weak equivalence}\inserted[typed]{aspheric} iff it is
``universally a weak equivalence'', namely iff for any base change $V\hand{\to}U$ in
$\hand{\mathcal{A}}$, the corresponding map $g\colon V'=V\mathrm{x}_U U'\hand{\to}V$ is a weak equivalence. M\over{e}{o}re-%
over, it is sufficient to check this property when V is a member of a
given generating family of $\hand{\mathcal{A}}$.

An interesting special case is the one when $\hand{\mathcal{A}}$ admits a family of
aspheric generators, or equivalently a \inserted[typed]{generating} full subcategory A whose objects
are aspheric in $\hand{\mathcal{A}}$ ; we will say in this case that the topos X (or $\hand{\mathcal{A}}$) is
\hand{*}\uline{locally aspheric} . This condition is satisfied for the most common topo-%
logical spaces (it suffices that each point admit a fundamental system of
contractible neighbourhoods), as well as for the topoi \inserted{$A\hat{}$} associated to
(essentially small) categories A (for such a topos, A itself is such a
\inserted[typed]{generating} full category made up with aspheric elements of $\hand{\mathcal{A}}$). In the\xout{\unsure{ir}} case\xout{s} of a
locally aspheric topos endowed with a generating full subcategory A as
above, a map $f\colon U'\hand{\longrightarrow} U$ is aspheric iff for any $a \hand{\in} \mathrm{Ob}\,A$ and map $a\hand{\longrightarrow} U$,
the object $U'\mathrm{x}_U a$ is aspheric (because for a map $\over{f}{g}\colon V'\hand{\to} V$ with V asphe-%
ric, g is a weak equivalence iff V' is aspheric\deleted{\illegible{1}} , which is a particular
case of the correponding statement valid for any map of topo\over{\unsure{s}}{i}).

Using this criterion, we get readily the following :
\uline{Proposition \hand{1}} Let $\hand{\mathcal{A}}$ be a locally aspheric topos, A a generating full sub-%
category made up with aspheric elements. The following conditions on $\hand{\mathcal{A}}$ are \margin{* loc. modelizing (cf § 39)}
equivalent\hand{:} \deleted{\illegible{1}}
a) Any aspheric element of $\hand{\mathcal{A}}$ is aspheric over e (the final object).
b) The product of any two aspheric elements of $\hand{\mathcal{A}}$ is again aspheric.
a'), b') as a), b) but with elements restricted to be in Ob A .

These conditions, in case $\hand{\mathcal{A}} = A\hand{\hat{}}$ are nothing but T 2) for A. \uline{It
turns out they imply already} T 1). M\over{e}{o}re generally, let's call a topos
